%!TEX TS-program = xelatex
%!TEX encoding = UTF-8 Unicode
% Curriculum Vitae - Dong-Joon Yi
% Built with Awesome-CV (https://github.com/posquit0/Awesome-CV)
% Build: ./build.sh   (or: xelatex cv.tex; xelatex cv.tex)

\documentclass[11pt, a4paper]{awesome-cv}

\geometry{left=1.5cm, top=1.0cm, right=1.5cm, bottom=1.6cm, footskip=.6cm}

% Allow long URLs/DOIs to break across lines, in the surrounding body font.
\usepackage{xurl}
\urlstyle{same}

% Underline that allows line breaks AND preserves adjacent spacing
% (standard \underline collapses surrounding elastic space in justified text).
% [normalem] keeps \emph as italics rather than underline.
\usepackage[normalem]{ulem}

% Reserve vertical space before a section so the heading isn't orphaned
% (used before Research Fellowship to keep heading + body on the same page).
\usepackage{needspace}

% Accent color (alternatives: awesome-emerald, awesome-skyblue, awesome-red,
% awesome-pink, awesome-orange, awesome-nephritis, awesome-concrete, awesome-darknight).
\colorlet{awesome}{awesome-skyblue}

\setbool{acvSectionColorHighlight}{true}
\renewcommand{\acvHeaderSocialSep}{\quad\textbar\quad}

%-------------------------------------------------------------------------------
%	STYLE OVERRIDES
%-------------------------------------------------------------------------------
% Source Sans 3's small-caps OpenType feature drops the dot on lowercase 'i',
% making "in" render as "iN" with a dotless i. Disable scshape everywhere it shows up.
% Also tighten the header name and bump small text for legibility.

\renewcommand*{\headerfirstnamestyle}[1]{{\fontsize{26pt}{1em}\headerfontlight\color{graytext} #1}}
\renewcommand*{\headerlastnamestyle}[1]{{\fontsize{26pt}{1em}\headerfont\bfseries\color{text} #1}}
\renewcommand*{\headerpositionstyle}[1]{{\fontsize{9pt}{1em}\bodyfont\itshape\color{awesome} #1}}

\renewcommand*{\entrypositionstyle}[1]{{\fontsize{9pt}{1em}\bodyfont\itshape\color{graytext} #1}}
\renewcommand*{\subentrypositionstyle}[1]{{\fontsize{8pt}{1em}\bodyfont\itshape\color{graytext} #1}}

% Subsection: bold awesome-color instead of small-caps.
\renewcommand*{\subsectionstyle}[1]{{\fontsize{11pt}{1em}\bodyfont\bfseries\color{awesome} #1}}

% Footer: drop scshape.
\renewcommand*{\footerstyle}[1]{{\fontsize{8pt}{1em}\footerfont\color{lighttext} #1}}

% Widen the cvhonor date column (default 1.5cm clips "2022, 2023") and the
% right-hand location column (2.5cm wraps long names like "Korean Physical Society").
\renewenvironment{cvhonors}{%
  \vspace{\acvSectionContentTopSkip}
  \vspace{-2mm}
  \begin{center}
    \setlength\tabcolsep{0pt}
    \setlength{\extrarowheight}{0pt}
    \begin{tabular*}{\textwidth}{@{\extracolsep{\fill}} L{2cm} L{\textwidth - 5.9cm} R{3.4cm}}
}{%
    \end{tabular*}
  \end{center}
}

%-------------------------------------------------------------------------------
%	PERSONAL INFORMATION
%-------------------------------------------------------------------------------
\name{Dong-Joon}{Yi}
\position{Ph.D. Candidate, Department of Electronic Engineering, Hanyang University}
\address{Semiconductor Nanophotonics Laboratory \, |\, Advisor: Prof. Mun Seok Jeong}

\email{iver.ydj@gmail.com}
\homepage{iverydj.github.io}
\googlescholar{MZE2mWsAAAAJ}{}
\extrainfo{iverydj@hanyang.ac.kr}

%-------------------------------------------------------------------------------
\begin{document}

\makecvheader

\makecvfooter
  {\today}
  {Dong-Joon Yi \, ·\, Curriculum Vitae}
  {\thepage}

%-------------------------------------------------------------------------------
%	EDUCATION & CAREERS
%-------------------------------------------------------------------------------
\cvsection{Education \& Careers}

\begin{cventries}

  \cventry
    {Ph.D. candidate in Electronic Engineering \, (GPA 4.3 / 4.5)}
    {Hanyang University}
    {Seoul, Republic of Korea}
    {Mar. 2024 -- Aug. 2027 (est.)}
    {
      \begin{cvitems}
        \item {Semiconductor Nanophotonics Laboratory}
        \item {Advisor: Prof. Mun Seok Jeong}
      \end{cvitems}
    }

  \cventry
    {M.S. in Solid-state Physics \, (GPA 4.5 / 4.5)}
    {Dongguk University}
    {Seoul, Republic of Korea}
    {Sep. 2022 -- Feb. 2024}
    {
      \begin{cvitems}
        \item {Advanced Physics on Electromaterial Laboratory}
        \item {Advisor: Prof. Kwun-Bum Chung}
        \item {Thesis: ``Performance Control of InSnZnO/InGaZnO Multi-Channel-Layered Transistors with Defect Analysis of Amorphous Oxide Semiconductors''}
      \end{cvitems}
    }

  \cventry
    {B.S. in Physics \, (GPA 4.2 / 4.5)}
    {Dongguk University}
    {Seoul, Republic of Korea}
    {Mar. 2016 -- Aug. 2022}
    {
      \begin{cvitems}
        \item {Military Service in the Republic of Korea (Jun. 2019 -- Apr. 2021)}
        \item {Undergraduate researcher in Advanced Physics on Electromaterial Laboratory (Jul. 2021 -- Aug. 2022)}
      \end{cvitems}
    }

\end{cventries}

%-------------------------------------------------------------------------------
%	SCHOLARSHIPS & AWARDS
%-------------------------------------------------------------------------------
% cvhonor signature: {bold-prominent}{gray-secondary}{location}{date}
% When 2nd arg is empty, no comma is inserted between fields 1 and 2.
\cvsection{Scholarships \& Awards}

\begin{cvhonors}
  \cvhonor{Physics Department Alumni Scholarship}{}{Dongguk University}{2022, 2023}
  \cvhonor{Dongguk 1.1.1 Scholarship}{}{Dongguk University}{2021}
  \cvhonor{Dongguk Future Talent Scholarship}{}{Dongguk University}{2022, 2023}
  \cvhonor{IC-PBL Graduate Championship}{3rd Place}{Hanyang University}{2024}
  \cvhonor{HY-BrainKorea G3 Award}{2nd Place}{Hanyang University}{2025}
  \cvhonor{Outstanding Presentation Award}{KPS Spring Meeting, Oral Session}{Korean Physical Society}{2026}
\end{cvhonors}

%-------------------------------------------------------------------------------
%	SKILLS
%-------------------------------------------------------------------------------
\cvsection{Skills}

\begin{cvskills}
  \cvskill
    {Spectroscopy}
    {Optical spectroscopy (Raman, Photoluminescence), Ellipsometry, Deep-Level Transient Spectroscopy (DLTS), Tip-Enhanced Nanospectroscopy}
  \cvskill
    {Machine Learning}
    {Physics-based data analysis, Deep learning for experimental data, Explainable Artificial Intelligence (XAI)}
  \cvskill
    {Simulation}
    {Finite-Difference Time-Domain (FDTD) simulations, Inverse design with adjoint optimization}
  \cvskill
    {Others}
    {Electrical characterization (I--V, etc.), X-ray Photoelectron Spectroscopy (XPS), Photolithography, Sputtering}
\end{cvskills}

%-------------------------------------------------------------------------------
%	RESEARCH
%-------------------------------------------------------------------------------
\cvsection{Research}

\begin{cvparagraph}
My research focuses on quantitative defect spectroscopy and experimental data analysis for two-dimensional (2D) materials and oxide semiconductors. I aim to extract physically meaningful parameters from electrical and optical measurements and relate them directly to material properties and device behavior, with an emphasis on interpretability and reproducibility.
\end{cvparagraph}

\cvsubsection{1.~~Defect \& Optical Spectroscopy}
\begin{cvitems}
  \item {Quantitative analysis of intrinsic and extrinsic defects in 2D materials and oxide semiconductors.}
  \item {Extraction of defect energy levels, densities, and carrier dynamics using Deep-Level Transient Spectroscopy (DLTS) and Photo-Induced Current Transient Spectroscopy (PICTS).}
  \item {Raman / Photoluminescence (PL) spectroscopy for probing defect states, doping, strain, and phase transitions.}
  \item {X-ray Photoelectron Spectroscopy (XPS) and spectroscopic ellipsometry (SE) for chemical composition, bonding states, and optical properties.}
  \item {Tip-Enhanced Raman Spectroscopy (TERS) for spatially resolved defect and structural inhomogeneity analysis beyond the diffraction limit.}
\end{cvitems}

\cvsubsection{2.~~Explainable Artificial Intelligence (XAI) for Experimental Data}
\begin{cvitems}
  \item {Development of analysis pipelines for spectroscopic and electrical datasets, including preprocessing, feature extraction, modeling, and validation.}
  \item {XAI to identify physically relevant data features.}
  \item {Connection of AI-extracted features to defect-related mechanisms and material properties under data-limited experimental conditions.}
\end{cvitems}

\cvsubsection{3.~~Advanced Electronic \& Photonic Devices}
\begin{cvitems}
  \item {Fabrication and characterization of electronic devices based on 2D materials and oxide semiconductors, including logic and neuromorphic devices.}
  \item {Quantitative analysis of device electrical responses (e.g., I--V characteristics and transient behavior) to identify dominant transport and trapping mechanisms.}
  \item {Finite-Difference Time-Domain (FDTD) simulations to evaluate optical responses and electric-field distributions in nano- to micro-scale structures, supporting the design and analysis of photonic components.}
\end{cvitems}

\begin{cvparagraph}
My work integrates defect-resolved spectroscopy, explainable analysis of experimental data, and device-level characterization to obtain quantitatively grounded insight into solid-state materials. The primary objective is to extract physically meaningful parameters from measurements and connect them directly to material behavior and device performance.
\end{cvparagraph}

%-------------------------------------------------------------------------------
%	PUBLICATIONS
%-------------------------------------------------------------------------------
\cvsection{Publications}

\begin{cvparagraph}
\href{https://scholar.google.co.kr/citations?user=MZE2mWsAAAAJ}{Google Scholar} \quad\textbar\quad \textsuperscript{\dag}\,First author \quad *\,Corresponding author
\end{cvparagraph}

% \me{} marks the CV owner in author lists (bold + underline). Uses ulem's
% \uline so that elastic spacing around the underlined hbox isn't collapsed.
% \pub{authors}{title}{venue}{year}{doi} appends one publication entry; the
% number is reverse-counted from \value{pubnum}. Bump pubnum by 1 when adding
% a new entry at the top. Same pattern for patents (patnum).
% \publink{authors}{title}{venue}{year}{full-url}{label} for entries that
% don't have a DOI (e.g., conference papers on OpenReview, arXiv).
\providecommand{\me}[1]{\textbf{\uline{#1}}}

\newcounter{pubnum}
\setcounter{pubnum}{11}  % (publication count) + 1; decrement each \pub
\providecommand{\pub}[5]{%
  \addtocounter{pubnum}{-1}%
  \item[\textbf{\arabic{pubnum}.}] {#1, ``#2,'' \textit{#3}, (#4), \href{https://doi.org/#5}{\nolinkurl{DOI:#5}}}%
}
\providecommand{\publink}[6]{%
  \addtocounter{pubnum}{-1}%
  \item[\textbf{\arabic{pubnum}.}] {#1, ``#2,'' \textit{#3}, (#4), \href{#5}{#6}}%
}

\begin{cvitems}
  \publink
    {Y.~Cho, J.~Yoo, S.~Yang, \me{D.-J. Yi}, S.~M.~Lee, M.~S.~Jeong, J.~Choo}
    {Position: Significant impact of numerical precision in scientific machine learning}
    {International Conference on Machine Learning (ICML), Position Paper Track}
    {2026}
    {https://openreview.net/forum?id=CbZbd4iorc}
    {OpenReview}

  \pub
    {H.~C.~Suh\textsuperscript{\dag}, T.~Kim\textsuperscript{\dag}, \me{D.-J. Yi}, D.~H.~Kim, S.~H.~Kim, J.~Yoo, S.~Bang, D.~Lee, J.~H.~Kim, Y.~S.~Won, K.~K.~Kim, M.~S.~Jeong*}
    {Prediction and assessment of nanoprobe for tip-enhanced Raman spectroscopy: Data-driven artificial intelligence approach}
    {Materials \& Design}
    {2026}
    {10.1016/j.matdes.2026.116066}

  \pub
    {T.~G. Weldemhret\textsuperscript{\dag}, \me{D.-J. Yi}\textsuperscript{\dag}, K.~Jeong, K.-B.~Chung*}
    {Evaluation of the high mobility and stability of InGaZnO/InSnZnO bilayer thin-film transistors via quantitative defect analysis}
    {Surfaces and Interfaces}
    {2025}
    {10.1016/j.surfin.2025.106870}

  \pub
    {S.~Bang\textsuperscript{\dag}, C.~Lee\textsuperscript{\dag}, D.~Choi, D.~Y.~Park, D.~H.~Kim, D.~Lee, \me{D.-J. Yi}, J.~Song, S.~J.~Yun, D.-W.~Kim*, M.~S.~Jeong*}
    {High Performance P-Channel Transistor Based on Amorphous Tellurium Trioxide}
    {Advanced Materials}
    {2025}
    {10.1002/adma.202504948}

  \pub
    {D.~H.~Kim\textsuperscript{\dag}, J.~Yoo, H.~C.~Suh, Y.~S.~Won, S.~H.~Kim, \me{D.-J. Yi}, B.~G.~Jeong, C.~Lee, D.~Lee, K.~K.~Kim, S.~M.~Lee, E.~K.~Koh, M.~S.~Jeong*}
    {Anomalous phonon softening with inherent strain in wrinkled monolayer WSe$_2$}
    {Advanced Materials}
    {2025}
    {10.1002/adma.202419414}

  \pub
    {K.~Jeong\textsuperscript{\dag}, D.~Y.~Shin, J.-M.~Park, \me{D.-J. Yi}, H.~Hong, H.-S.~Kim*, K.-B.~Chung*}
    {Noncontact Monitoring and Imaging of the Operation and Performance of Thin-Film Field-Effect Transistors}
    {Advanced Science}
    {2025}
    {10.1002/advs.202407923}

  \pub
    {A.~Song\textsuperscript{\dag}, M.~J.~Kim, \me{D.-J. Yi}, S.~Kwon, D.-W.~Kim, S.~Kim, J.-H.~Bae, S.~Park, Y.~S.~Rim*, K.-S.~Jeong*, K.-B.~Chung*}
    {Control of sensitivity in metal oxide electrolyte gated field-effect transistor-based glucose sensor by electronegativity modulation}
    {Scientific Reports}
    {2024}
    {10.1038/s41598-024-76885-x}

  \pub
    {H.~Hong\textsuperscript{\dag}, M.~J.~Kim, \me{D.-J. Yi}, D.~Y.~Shin, Y.-K.~Moon, K.-B.~Chung*}
    {Quantitative dynamic evolution of unoccupied states in hydrogen diffused InGaZnSnO TFT under positive bias temperature stress}
    {ACS Applied Electronic Materials}
    {2024}
    {10.1021/acsaelm.4c01430}

  \pub
    {H.~Hong\textsuperscript{\dag}, \me{D.-J. Yi}, Y.-K.~Moon, K.-S.~Son, J.~H.~Lim, K.~Jeong, K.-B.~Chung*}
    {Double gated a-InGaZnO TFT properties based on quantitative defect analysis and computational modeling}
    {IEEE Transactions on Electron Devices}
    {2024}
    {10.1109/TED.2023.3347503}

  \pub
    {H.~Hong\textsuperscript{\dag}, M.~J.~Kim, \me{D.-J. Yi}, Y.-K.~Moon, K.-S.~Son, J.~H.~Lim, K.~Jeong*, K.-B.~Chung*}
    {Quantitative analysis of defect states in InGaZnO within 2 eV below the conduction band via photo-induced current transient spectroscopy}
    {Scientific Reports}
    {2023}
    {10.1038/s41598-023-40162-0}
\end{cvitems}

%-------------------------------------------------------------------------------
%	PATENTS
%-------------------------------------------------------------------------------
\cvsection{Patents}

\newcounter{patnum}
\setcounter{patnum}{4}  % (patent count) + 1; decrement each \patentitem
\providecommand{\patentitem}[1]{%
  \addtocounter{patnum}{-1}%
  \item[\textbf{\arabic{patnum}.}] {#1}%
}

\begin{cvitems}
  \patentitem
    {K.-B.~Chung, K.~Jeong, \me{D.-J. Yi}, (Dongguk University Industry Academy Cooperation Foundation). ``Device Array Analysis Method and Analysis Apparatus Therefor.'' \textit{US Patent}, 18771046.}
  \patentitem
    {K.-B.~Chung, K.~Jeong, \me{D.-J. Yi}, (Dongguk University Industry Academy Cooperation Foundation). ``Analysis Method of Semiconductor Device and Analysis Apparatus for the Method.'' \textit{Patent, Republic of Korea}, 10-2024-0049532.}
  \patentitem
    {K.-B.~Chung, K.~Jeong, \me{D.-J. Yi}, H.~Hong, (Dongguk University Industry Academy Cooperation Foundation). ``Method of Confirming Performance of Secondary Battery and Apparatus of Performing the Same.'' \textit{Patent, Republic of Korea}, 10-2023-0161220.}
\end{cvitems}

%-------------------------------------------------------------------------------
%	RESEARCH FELLOWSHIP
%-------------------------------------------------------------------------------
% Keep heading + body together on the same page (otherwise the heading orphans
% at the bottom of the previous page while the entry overflows to a new page).
\needspace{6\baselineskip}
\cvsection{Research Fellowship}

\begin{cventries}
  \cventry
    {Principal Investigator}
    {NRF Basic Science Research Program}
    {50\,M KRW}
    {Sep. 2025 -- Aug. 2027}
    {
      \begin{cvitems}
        \item {Project: ``Deep-Level Transient Spectroscopy and Nanospectroscopy for Ultrafine Defect Analysis Using Explainable AI in 2D Quantum Semiconductors''}
        \item {Funded by the Ministry of Education through the National Research Foundation of Korea (NRF).}
      \end{cvitems}
    }
\end{cventries}

\end{document}
